\section{问题分析}
\subsection{问题分析}
先从整题角度说明各问怎样构成一条研究主线，区分共用资料与每问独有输入；下文再逐问解释本问困难、数据条件及方法思路。

\questionlabel{针对问题一，}填写研究目标、可直接使用的资料、需要构造的变量，以及为什么选择相应方法。

\questionlabel{针对问题二，}填写对前问信息的实际依赖，并说明新增条件怎样改变模型或求解目标。

\questionlabel{针对问题三，}填写需要进一步解决的疑点、候选路线和相应验证依据。

\questionlabel{针对问题四，}填写本问的建模与检验思路，以及如何综合前问有效成果。

需要用方法图解释整题关系时，可替换图~\ref{fig:roadmap}的占位框；无此需要时删除整个图环境，并同步正文引用。
\begin{figure}[htbp]
\centering
\fbox{\parbox[c][18mm][c]{0.88\textwidth}{\centering 图形位置示例：研究对象、方法与输出之间的关系\\此框不是计算结果或正式技术路线图}}
\caption{整题研究技术路线（格式位置示例）}\label{fig:roadmap}
\end{figure}
\FloatBarrier
